\chapter{Alkohol: Mengaktifkan Gugus OH}\label{ch:b1:alcohol-activation}

Alkohol ada di mana-mana: dibuat berton-ton dari alkena dan dari
fermentasi, dihasilkan oleh setiap adisi Grignard pada bab sebelumnya.
Namun, alkohol tidak mengalami substitusi seperti haloalkana: gugus OH
harus pergi sebagai ion hidroksida, sebuah \omterm{def:b1:predominant-reaction:strong-acid}{basa kuat} dan \omterm{def:b1:electronic-effects:nucleophile}{gugus pergi} yang
sangat buruk (\cref{ch:b1:nucleophilic-substitution}). Para kimiawan
mempunyai tiga jawaban. Lepaskan proton OH-nya, maka \omterm{def:b1:alcohol-activation:alkoxide}{alkoksidanya} menjadi
\omterm{def:b1:electronic-effects:nucleophile}{nukleofil} yang kuat. Tambahkan sebuah proton, maka \omterm{def:b1:electronic-effects:nucleophile}{gugus perginya} menjadi
air. Atau gantikan hidrogennya dengan gugus yang menjadikan oksigen itu
bagian dari \omterm{def:b1:electronic-effects:nucleophile}{gugus pergi} yang sangat baik, tanpa menyentuh karbonnya. Bab
ini mengikuti ketiga jalan itu, beserta pilihan mekanisme dan
stereokimia yang menyertainya.

\begin{recall}
SN1, SN2, dan stereokimianya (\cref{ch:b1:nucleophilic-substitution});
E1, E2, dan \omterm{prop:b1:elimination:zaitsev}{aturan Zaitsev} (\cref{ch:b1:elimination}); \omterm{def:b1:predominant-reaction:acidity-constant}{tetapan keasaman}
dan metode \omterm{def:b1:predominant-reaction:predominant-reaction}{reaksi predominan} (\cref{ch:b1:predominant-reaction});
pengaruh gugus alkil terhadap keasaman alkohol
(\cref{exo:b1:electronic-effects:11}).
\end{recall}

\section{Alkohol sebagai asam: alkoksida}

\begin{definition}[Alkoksida]\label{def:b1:alcohol-activation:alkoxide}
\emph{Alkoksida}\index{alkoksida} adalah basa konjugasi \ce{RO-} dari
alkohol \ce{ROH}: metoksida \ce{CH3O-}, etoksida \ce{CH3CH2O-},
tert-butoksida \ce{(CH3)3CO-}. Alkoksida adalah \omterm{def:b1:predominant-reaction:strong-acid}{basa kuat} dan \omterm{def:b1:electronic-effects:nucleophile}{nukleofil}
yang baik.
\end{definition}

\begin{proposition}[Keasaman alkohol]\label{prop:b1:alcohol-activation:alcohol-acidity}
Alkohol adalah asam yang lebih lemah daripada air atau kira-kira sama
lemahnya: $\mathrm{p}K_a$ 15.5 (metanol), 15.9 (etanol), 17.1
(propan-2-ol), 19.2 (2-metilpropan-2-ol), dibandingkan dengan 14.0 untuk
pasangan \ce{H2O}/\ce{OH-}. Karena itu hidroksida hanya mendeprotonasinya
sebagian; konversi sempurna menjadi \omterm{def:b1:alcohol-activation:alkoxide}{alkoksida} memerlukan basa yang lebih
kuat atau logam alkali.
\end{proposition}
% ledger: pka:methanol, pka:ethanol, pka:2-propanol, pka:tBuOH

\begin{proof}
Reaksi \ce{ROH + OH- <=> RO- + H2O} mempunyai $K = 10^{14.0 -
\mathrm{p}K_a(\ce{ROH})}$, antara $10^{-1.5}$ dan $10^{-5}$: tidak
menguntungkan. Logam natrium mereduksi protonnya,
\ce{2ROH + 2Na -> 2RO- + 2Na+ + H2}, dan ion hidrida dari natrium hidrida
mengambilnya secara tak reversibel, \ce{ROH + H- -> RO- + H2}: hidrogennya
lepas, dan reaksinya sempurna berapa pun tetapan kesetimbangannya.
\end{proof}

\begin{inthelab}[Natrium hidrida]
Natrium hidrida dijual sebagai serbuk abu-abu yang terdispersi dalam
minyak mineral, yang melindunginya dari kelembapan udara. Zat itu bereaksi
hebat dengan air dan melepaskan hidrogen yang mudah menyala; zat itu
ditimbang dengan cepat, ditambahkan sedikit demi sedikit ke dalam alkohol
yang kering di bawah gas inert, dan gelembung hidrogen memperlihatkan
kemajuan deprotonasinya. Hidrida yang berlebih dirusak pada akhirnya
dengan menambahkan sebuah alkohol secara perlahan, tidak pernah air.
\end{inthelab}

\section{Sintesis eter Williamson}

\begin{definition}[Sintesis eter Williamson]\label{def:b1:alcohol-activation:williamson}
\emph{Sintesis eter Williamson}\index{sintesis eter Williamson}
membuat eter $\mathrm{R{-}O{-}R'}$ melalui reaksi SN2 sebuah \omterm{def:b1:alcohol-activation:alkoxide}{alkoksida} \ce{RO-}
pada haloalkana (atau \omterm{def:b1:alcohol-activation:sulfonate}{ester sulfonat}) $\mathrm{R'{-}X}$.
\end{definition}

\begin{omfigure}
\setchemfig{atom sep=2.1em}
\schemestart
\chemfig{CH_3CH_2-@{o}O^{-}}\,\chemfig{Na^{+}}
\+
\chemfig{@{c}CH_3-[@{ci}]@{i}I}
\arrow{->}
\chemfig{CH_3CH_2-O-CH_3}
\+
\chemfig{Na^{+}}\,\chemfig{I^{-}}
\schemestop

\medskip
\chemmove{\draw[->, thick, shorten <=2pt, shorten >=2pt](o) .. controls +(80:8mm) and +(100:8mm) .. (c);
          \draw[->, thick, shorten <=2pt, shorten >=2pt](ci) .. controls +(90:5mm) and +(90:5mm) .. (i);}

{\small Sintesis Williamson metoksietana: etoksida menyerang karbon
iodometana dari sisi yang berseberangan dengan iodin (SN2).}
\end{omfigure}

\begin{method}[Memilih kedua pasangan reaksi]\label{met:b1:alcohol-activation:williamson-choice}
Eter $\mathrm{R{-}O{-}R'}$ dapat dipotong di kedua sisi oksigennya.
\begin{enumerate}
  \item Tuliskan kedua pasangan yang mungkin: \ce{RO-} dengan $\mathrm{R'X}$, dan
        $\mathrm{R'O^-}$ dengan \ce{RX}.
  \item Pertahankan pasangan yang halidanya metil atau primer: SN2
        memerlukan karbon yang tidak terhalang.
  \item Tolak pasangan dengan halida tersier: \omterm{def:b1:alcohol-activation:alkoxide}{alkoksida}, sebuah \omterm{def:b1:predominant-reaction:strong-acid}{basa kuat},
        malah akan melakukan eliminasi (E2).
  \item Dengan halida sekunder, harapkan campuran eter dan alkena.
\end{enumerate}
\end{method}

\begin{example}[Eter yang tidak simetris]\label{ex:b1:alcohol-activation:williamson}
2-Metoksi-2-metilpropana (metil tert-butil eter): natrium tert-butoksida
dengan iodometana menghasilkannya dengan bersih; natrium metoksida dengan
bromida tersier, 2-bromo-2-metil\-propana, menghasilkan 2-metilpropena
melalui E2. Gugus meruah di sisi \omterm{def:b1:alcohol-activation:alkoxide}{alkoksida}, gugus kecil di sisi halida.
\end{example}

\section{Mengaktifkan gugus OH}

\begin{definition}[Ester sulfonat]\label{def:b1:alcohol-activation:sulfonate}
\emph{Ester sulfonat}\index{ester sulfonat} $\mathrm{R{-}O{-}SO_2{-}R'}$ terbentuk
dari alkohol dan sulfonil klorida $\mathrm{R'SO_2Cl}$ dengan adanya
basa (piridina). \emph{Tosilat}\index{tosilat}, \ce{R-OTs}, berasal
dari 4-metilbenzenasulfonil klorida (tosil klorida);
\emph{mesilat}\index{mesilat}, \ce{R-OMs}, dari metanasulfonil klorida.
Anion sulfonat $\mathrm{R'SO_3^-}$, basa konjugasi sebuah \omterm{def:b1:predominant-reaction:strong-acid}{asam kuat} yang
distabilkan oleh resonansi atas tiga oksigen, adalah \omterm{def:b1:electronic-effects:nucleophile}{gugus pergi} yang sangat baik.
\end{definition}

\begin{omfigure}
\setchemfig{atom sep=2em}
\schemestart
\chemfig{R-O-H}
\+
\chemfig{Cl-S(=[2]O)(=[6]O)-[:0]*6(-=-(-CH_3)=-=)}
\arrow{->[piridina]}
\chemfig{R-O-S(=[2]O)(=[6]O)-[:0]*6(-=-(-CH_3)=-=)}
\schemestop
\par\vspace{6pt}

{\small Tosilasi sebuah alkohol. Ikatan O--H digantikan oleh O--S; ikatan
C--O alkohol tidak tersentuh. Piridina mengikat HCl yang terbentuk.}
\end{omfigure}

\begin{proposition}[Konfigurasi melalui tosilasi dan substitusi]\label{prop:b1:alcohol-activation:tosylation-retention}
Tosilasi alkohol pada karbon stereogenik mempertahankan \omterm{def:b1:stereochemistry-in-depth:isomers}{konfigurasinya}
(retensi); substitusi SN2 berikutnya terhadap \omterm{def:b1:alcohol-activation:sulfonate}{tosilat} membaliknya. Kedua
tahap bersama-sama menggantikan OH dengan sebuah \omterm{def:b1:electronic-effects:nucleophile}{nukleofil} disertai inversi.
\end{proposition}

\begin{proof}
Dalam tosilasi, ikatan yang berubah adalah O--H dan S--Cl: oksigen alkohol
menyerang belerang dan hidrogennya berpindah ke basa. Tidak ada ikatan ke
karbon stereogenik yang putus, sehingga susunan di sekitarnya tetap sama.
Dalam tahap SN2, ikatan C--O putus sementara \omterm{def:b1:electronic-effects:nucleophile}{nukleofil} berikatan di sisi
yang berseberangan (\cref{prop:b1:nucleophilic-substitution:sn2-inversion}).
\end{proof}

\begin{omfigure}
\begin{tikzpicture}[font=\small, >={Stealth}]
  \node (a) at (0,0) {\chemfig{C(-[:150]H_3C)(<[:210]H)(<:[:250]C_2H_5)-OH}};
  \node at (0,-1.4) {\cip{S}-butan-2-ol};
  \draw[->, thick] (1.3,0) -- (2.5,0) node[midway, above] {TsCl} node[midway, below, font=\scriptsize] {retensi};
  \node (b) at (3.9,0) {\chemfig{C(-[:150]H_3C)(<[:210]H)(<:[:250]C_2H_5)-OTs}};
  \node at (3.9,-1.4) {tosilat \cip{S}};
  \draw[->, thick] (5.3,0) -- (6.6,0) node[midway, above] {\ce{CN-}} node[midway, below, font=\scriptsize] {SN2, inversi};
  \node (c) at (8.3,0) {\chemfig{[:0]NC-C(-[:30]CH_3)(<[:-30]H)(<:[:-70]C_2H_5)}};
  \node at (8.3,-1.4) {nitril \cip{R}};
\end{tikzpicture}

{\small Dari \cip{S}-butan-2-ol ke 2-metilbutananitril: tosilasi
mempertahankan \omterm{def:b1:stereochemistry-in-depth:isomers}{konfigurasinya}, tahap SN2 dengan sianida membaliknya.
Secara keseluruhan: OH digantikan oleh CN disertai inversi.}
\end{omfigure}

Tujuan yang sama dapat dicapai dalam asam: gugus OH yang terprotonasi
pergi sebagai air, \ce{R-OH2+ -> R+ + H2O} (SN1, E1) atau karena serangan
(SN2). Akan tetapi, asam membatasi \omterm{def:b1:electronic-effects:nucleophile}{nukleofilnya} pada yang tahan terhadapnya
--- ion halida, air, alkohol --- dan membuka pintu bagi \omterm{def:b1:electronic-effects:intermediates}{karbokation}
beserta reaksi sampingnya.

\section{Konversi menjadi haloalkana}

\begin{proposition}[Alkohol dan hidrogen halida]\label{prop:b1:alcohol-activation:hx-by-class}
Alkohol tersier menghasilkan haloalkana dengan asam klorida pekat pada
suhu kamar (SN1); alkohol primer memerlukan asam bromida atau asam iodida
dan pemanasan (SN2 pada alkohol yang terprotonasi); alkohol sekunder
bereaksi dengan laju pertengahan, sering melalui kedua mekanisme.
\end{proposition}

\begin{proof}
Protonasi mengubah OH menjadi \ce{-OH2+}. Untuk alkohol tersier, air
mudah pergi dan menghasilkan \omterm{def:b1:electronic-effects:intermediates}{karbokation} stabil yang ditangkap oleh
halida (SN1): \omterm{def:b1:electronic-effects:nucleophile}{nukleofil} sesederhana \ce{Cl-} pun cukup. \omterm{def:b1:electronic-effects:intermediates}{Karbokation} primer
terlalu tidak stabil; halida harus mendorong air keluar dari belakang
(SN2), yang memerlukan \omterm{def:b1:electronic-effects:nucleophile}{nukleofil} yang baik (\ce{Br-}, \ce{I-}, lebih baik
daripada \ce{Cl-} dalam medium protik,
\cref{prop:b1:electronic-effects:nucleophilicity-basicity}) dan kalor.
\end{proof}

\begin{example}[Uji penggolongan]\label{ex:b1:alcohol-activation:lucas}
Campuran asam klorida pekat dan seng klorida (\omterm{def:b1:lewis-resonance-vsepr:lewis-acid}{asam Lewis} yang membantu
OH pergi) membuat alkohol tersier keruh seketika --- klorida yang tidak
larut memisah ---, alkohol sekunder dalam beberapa menit, dan alkohol
primer hampir tidak sama sekali pada suhu kamar: urutan kestabilan
\omterm{def:b1:electronic-effects:intermediates}{karbokation} yang terlihat di dalam tabung reaksi.
\end{example}

\section{Dehidrasi dan pembentukan eter dalam asam}

\begin{definition}[Dehidrasi alkohol]\label{def:b1:alcohol-activation:dehydration}
\emph{Dehidrasi}\index{dehidrasi} alkohol adalah eliminasi air, yang
dikatalisis oleh \omterm{def:b1:predominant-reaction:strong-acid}{asam kuat} (asam sulfat atau asam fosfat) dan kalor,
menghasilkan alkena: $\mathrm{R_2CH{-}CR'_2OH} \to \mathrm{R_2C{=}CR'_2} + \ce{H2O}$.
\end{definition}

\begin{omfigure}
\setchemfig{atom sep=2.1em}
\schemestart
\chemfig{C(-[2]CH_3)(-[:210]H_3C)(-[:-30]CH_2CH_3)-OH}
\arrow{->[\ce{H+}]}
\chemfig{C(-[2]CH_3)(-[:210]H_3C)(-[:-30]CH_2CH_3)-OH_2^{+}}
\arrow{->[\ce{-H2O}]}
\chemfig{C^{+}(-[2]CH_3)(-[:210]H_3C)(-[:-30]CH_2CH_3)}
\schemestop

\medskip
\schemestart
\chemfig{C^{+}(-[2]CH_3)(-[:210]H_3C)(-[:-30]CH_2CH_3)}
\arrow{->[\ce{-H+}][(utama)]}
\chemfig{C(-[2]CH_3)(-[:210]H_3C)=[:-30]CHCH_3}
\schemestop

\medskip
{\small \omterm{def:b1:alcohol-activation:dehydration}{Dehidrasi} 2-metilbutan-2-ol yang dikatalisis asam (E1): protonasi,
lepasnya air menjadi \omterm{def:b1:electronic-effects:intermediates}{karbokation} tersier, lepasnya sebuah proton.
Melepaskan proton dari \ce{CH2} menghasilkan 2-metilbut-2-ena yang
trisubstitusi (Zaitsev, utama); dari gugus metil, 2-metilbut-1-ena (minor).}
\end{omfigure}

\begin{proposition}[Eter atau alkena]\label{prop:b1:alcohol-activation:ether-vs-alkene}
Alkohol primer yang dipanaskan dengan asam sulfat terutama menghasilkan
eter simetris pada suhu sedang (satu molekul menyubstitusi molekul lain
melalui SN2) dan terutama alkena pada suhu yang lebih tinggi; alkohol
sekunder dan tersier menghasilkan alkena.
\end{proposition}

\begin{proof}
Untuk alkohol primer, alkohol yang terprotonasi dapat diserang pada
karbonnya oleh molekul alkohol kedua (SN2, menghasilkan \ce{R-O-R} sesudah
lepasnya sebuah proton) atau pada hidrogen $\beta$ (E2). Eliminasi
menghasilkan lebih banyak molekul daripada yang dipakainya dan makin
disukai ketika suhu naik (\cref{prop:b1:elimination:sn-vs-e}); substitusi
mendominasi ketika suhunya lebih rendah. Alkohol sekunder dan tersier
terionisasi menjadi \omterm{def:b1:electronic-effects:intermediates}{karbokation}, yang lebih mudah melepaskan proton
daripada ditangkap oleh \omterm{def:b1:electronic-effects:nucleophile}{nukleofil} lemah, terlebih lagi pada pemanasan;
air yang terbentuk juga dikeluarkan dengan mendistilasi alkena yang
mudah menguap.
\end{proof}

\section{Latihan}

\begin{exercise}[$\star$]\label{exo:b1:alcohol-activation:1}
Tuliskan reaksi etanol dengan natrium, dengan natrium hidrida, dan dengan
natrium hidroksida; hitunglah tetapan reaksi yang terakhir ($\mathrm{p}K_a$
etanol 15.9).
\end{exercise}

\begin{exercise}[$\star$]\label{exo:b1:alcohol-activation:2}
Tentukan produk dari: natrium metoksida dengan 1-bromopropana; natrium
etoksida dengan iodometana; natrium fenoksida dengan bromoetana.
\end{exercise}

\begin{exercise}[$\star$]\label{exo:b1:alcohol-activation:3}
Tuliskan tosilasi propan-1-ol dan reaksi \omterm{def:b1:alcohol-activation:sulfonate}{tosilatnya} dengan natrium
iodida. Apa keuntungannya dibandingkan reaksi langsung alkohol itu dengan
natrium iodida?
\end{exercise}

\begin{exercise}[$\star$]\label{exo:b1:alcohol-activation:4}
Tentukan alkena-alkena yang terbentuk pada \omterm{def:b1:alcohol-activation:dehydration}{dehidrasi} asam butan-2-ol dan
2-metilpentan-2-ol, serta alkena utamanya.
\end{exercise}

\begin{exercise}[$\star\star$]\label{exo:b1:alcohol-activation:5}
Usulkan sintesis Williamson untuk 1-etoksi-2-metilpropana dan jelaskan
pilihan pasangan reaksinya. Mengapa pilihan yang lain akan menghasilkan alkena?
\end{exercise}

\begin{exercise}[$\star\star$]\label{exo:b1:alcohol-activation:6}
\cip{R}-Oktan-2-ol diubah menjadi \omterm{def:b1:alcohol-activation:sulfonate}{tosilatnya}, lalu direaksikan dengan
natrium etoksida dalam etanol. Tentukan \omterm{def:b1:stereochemistry-in-depth:isomers}{konfigurasi} eternya. Apa yang
akan diperoleh dengan memanaskan alkohol yang sama dengan asam sulfat?
\end{exercise}

\begin{exercise}[$\star\star$]\label{exo:b1:alcohol-activation:7}
Tuliskan \omterm{def:b1:elementary-steps:elementary-step}{mekanisme reaksi} 2-metilpropan-2-ol dengan asam klorida pekat
dan jelaskan mengapa reaksinya cepat pada suhu kamar.
\end{exercise}

\begin{exercise}[$\star\star$]\label{exo:b1:alcohol-activation:8}
Tuliskan mekanisme pembentukan etoksietana dari etanol dalam asam sulfat,
beserta eliminasi yang bersaing dengannya.
\end{exercise}

\begin{exercise}[$\star\star$]\label{exo:b1:alcohol-activation:9}
Dalam uji penggolongan pada bab ini, alkohol manakah di antara butan-1-ol,
butan-2-ol, dan 2-metilpropan-2-ol yang paling dahulu menjadi keruh?
Tuliskan produknya.
\end{exercise}

\begin{exercise}[$\star\star\star$]\label{exo:b1:alcohol-activation:10}
3,3-Dimetilbutan-2-ol, yang dipanaskan dalam asam, terutama menghasilkan
2,3-dimetilbut-2-ena, yang kerangka karbonnya berbeda dari kerangka
alkoholnya. Usulkan mekanisme dengan pergeseran-1,2 sebuah gugus metil
dalam \omterm{def:b1:electronic-effects:intermediates}{karbokation} itu, dan berikan alasan pergeseran itu.
\end{exercise}

\begin{exercise}[$\star\star\star$]\label{exo:b1:alcohol-activation:11}
Rancanglah sintesis \cip{S}-2-metoksibutana dari \cip{R}-butan-2-ol, dan
sintesis lain dari \cip{S}-butan-2-ol, masing-masing dengan \omterm{def:b1:stereochemistry-in-depth:isomers}{konfigurasi}
yang tetap terkendali.
\end{exercise}

\begin{exercise}[$\star\star\star$]\label{exo:b1:alcohol-activation:12}
Tunjukkan bahwa kesetimbangan \ce{2CH3CH2OH <=> (CH3CH2)2O + H2O} dapat
digeser ke arah eter dengan mengeluarkan air secara terus-menerus, dengan
memakai \omterm{def:b1:extent-q-and-k:quotient}{kuosien reaksi} (\cref{ch:b1:extent-q-and-k}).
\end{exercise}

\section{Soal: Dua Jalan Menuju Sebuah Eter}

\begin{problem}[{Soal akhir pekan --- dua jalan Williamson menuju metil
tert-butil eter, eter yang stereokimianya terkendali dari alkohol yang
aktif optis, dehidrasi alkohol tersier, dan massa eter yang
diperoleh}]
\label{pb:b1:alcohol-activation:1}
Pada masa lalu, 2-metoksi-2-metilpropana (metil tert-butil eter, MTBE)
diproduksi dalam skala besar sebagai aditif bahan bakar. Massa molar
(\unit{g/mol}): \ce{C4H10O} 74.0, \ce{C5H12O} 88.0, \ce{CH3I} 141.9.

\medskip
\noindent\textbf{Bagian I --- MTBE melalui sintesis Williamson.}
\begin{enumerate}
  \item Tuliskan kedua pasangan reaksi (\omterm{def:b1:alcohol-activation:alkoxide}{alkoksida} dan halida) yang dapat
        menghasilkan MTBE.
  \item Pasangan manakah yang gagal? Tuliskan reaksi yang terjadi sebagai
        gantinya.
  \item Tuliskan reaksi yang berhasil beserta mekanismenya.
  \item Bagaimana natrium tert-butoksida dibuat dari 2-metilpropan-2-ol?
        Mengapa tidak dengan natrium hidroksida?
  \item Mengapa reaksinya dijalankan dalam alkohol yang kering atau dalam
        \omterm{def:b1:intermolecular-forces-solvents:solvent-classes}{pelarut aprotik}?
  \item Apakah produknya \omterm{def:b1:stereochemistry-in-depth:stereocentre}{kiral}?
\end{enumerate}

\medskip
\noindent\textbf{Bagian II --- Eter yang aktif optis.}
\begin{enumerate}[resume]
  \item \cip{R}-Butan-2-ol diubah menjadi \omterm{def:b1:alcohol-activation:sulfonate}{tosilatnya}. Tuliskan reaksinya
        dan tentukan \omterm{def:b1:stereochemistry-in-depth:isomers}{konfigurasi} \omterm{def:b1:alcohol-activation:sulfonate}{tosilat} itu.
  \item \omterm{def:b1:alcohol-activation:sulfonate}{Tosilat} itu bereaksi dengan natrium metoksida. Tuliskan produknya
        beserta \omterm{def:b1:stereochemistry-in-depth:isomers}{konfigurasinya}.
  \item \omterm{def:b1:stereochemistry-in-depth:isomers}{Konfigurasi} alkohol manakah yang menghasilkan
        \cip{R}-2-metoksibutana melalui jalan yang sama?
  \item Jalan lain: \omterm{def:b1:alcohol-activation:alkoxide}{alkoksida} \cip{R}-butan-2-ol dengan iodometana.
        Bagaimana \omterm{def:b1:stereochemistry-in-depth:isomers}{konfigurasi} produknya? Mengapa?
  \item Reaksi samping apakah yang bersaing pada pertanyaan 8, dan mengapa
        reaksi itu terbatas di sini?
  \item Mengapa pemanasan \cip{R}-butan-2-ol dengan metanol dalam asam
        sulfat akan menghasilkan eter yang hampir rasemat, kalaupun eter
        terbentuk?
\end{enumerate}

\medskip
\noindent\textbf{Bagian III --- \omterm{def:b1:alcohol-activation:dehydration}{Dehidrasi} alkohol tersier.}
\begin{enumerate}[resume]
  \item Tuliskan mekanisme \omterm{def:b1:alcohol-activation:dehydration}{dehidrasi} asam 2-metilbutan-2-ol.
  \item Tentukan kedua alkenanya dan ramalkan alkena utamanya.
  \item Mengapa alkohol ini tidak membentuk eter?
  \item Mengapa 2-metilbut-2-ena didistilasi keluar segera setelah
        terbentuk?
  \item Gambarlah \omterm{def:b1:elementary-steps:energy-profile}{profil energi} tahap E1 dan tandai \omterm{def:b1:elementary-steps:rds}{tahap penentu lajunya}.
  \item Dapatkah alkohol primer didehidrasi melalui mekanisme yang sama?
\end{enumerate}

\medskip
\noindent\textbf{Bagian IV --- Massa.}
\begin{enumerate}[resume]
  \item \qty{10.0}{g} 2-metilpropan-2-ol diubah menjadi \omterm{def:b1:alcohol-activation:alkoxide}{alkoksidanya}.
        Hitunglah jumlahnya.
  \item Berapa massa iodometana yang diperlukan untuk kelebihan
        \qty{10}{\percent}?
  \item Hitunglah massa teoretis MTBE.
  \item Hitunglah massa yang diperoleh untuk \omterm{def:b1:extent-q-and-k:fractional-extent}{rendemen} \qty{75}{\percent}.
\end{enumerate}
\end{problem}
